% Pandoc (LuaLaTeX) já carrega fontspec e babel[brazilian] via lang: pt-BR no frontmatter.
% Este arquivo adiciona apenas o que o template padrão não configura.

% Geometria A4 com margens confortáveis
\usepackage[a4paper, top=2.5cm, bottom=2.5cm, left=2.5cm, right=2.5cm]{geometry}

% Quebra automática de linhas em blocos de código
% "Highlighting" é o ambiente que Pandoc usa com syntax highlighting
\usepackage{fvextra}
\DefineVerbatimEnvironment{Highlighting}{Verbatim}{%
  breaklines=true,%
  breakanywhere=true,%
  breaksymbol={},%
  commandchars=\\\{\}%
}

% Cores (necessário antes de hyperref)
\usepackage{xcolor}

% Tabelas — wrap automático de texto nas células
\usepackage{longtable}
\usepackage{booktabs}
\usepackage{array}
\newcolumntype{L}[1]{>{\raggedright\arraybackslash}p{#1}}

% Cabeçalho e rodapé personalizados
\usepackage{fancyhdr}
\pagestyle{fancy}
\fancyhf{}
\fancyhead[L]{\small\textit{WBA Windows Toolkit}}
\fancyhead[R]{\small\textit{Manual do Operador}}
\fancyfoot[C]{\thepage}
\renewcommand{\headrulewidth}{0.4pt}
\renewcommand{\footrulewidth}{0pt}

% Fonte monoespaçada para blocos de código (fontspec já foi carregado pelo Pandoc)
\AtBeginDocument{%
  \IfFontExistsTF{Libertinus Mono}{%
    \setmonofont{Libertinus Mono}[Scale=0.9]%
  }{%
    \setmonofont{Latin Modern Mono}[Scale=0.9]%
  }%
}
